Automated tools can detect potential vulnerabilities at scale: taint analyzers flag suspicious data flows, fuzzers produce crashes, LLM scanners identify dangerous patterns. But detection at scale produces thousands of candidates, and a candidate is not a confirmed vulnerability. The gap between ``this might be a bug'' and ``this is a real, exploitable vulnerability'' is where automated discovery breaks down. This dissertation presents 3 approaches that progressively close that gap through stronger forms of validation evidence.

First, Operation Mango (USENIX Security 2024) bakes partial validation into the detection phase itself. MangoDFA, a novel static data-flow analysis for firmware binaries, tracks value transformation history through Rich Expressions rather than binary taint bits, producing candidates that carry evidence about whether each value reaching a vulnerability sink is still exploitable. On the standard Karonte dataset, Mango analyzed all 3,599 binaries (a 27$\times$ increase in scope over prior work) while averaging 8 minutes per binary. On a large-scale corpus of 1,698 firmware images containing 770,374 binaries, Mango produced 10,834 ranked vulnerability alerts and confirmed 166 exploitable vulnerabilities.

Second, Artiphishell (DEF CON 2025) validates through multi-technique convergence. Built for the DARPA AI Cyber Challenge, Artiphishell is a fully autonomous Cyber Reasoning System that coordinates coverage-guided fuzzing, grammar-based fuzzing, LLM-powered scanning, root-cause analysis, and AI-driven patching through a shared Analysis Graph. When multiple independent techniques point to the same bug and a patch fixes it, the system has built a convergent case for the vulnerability. Artiphishell advanced to the AIxCC semi-finals.

Third, we propose a symbolic execution engine operating on angr's AIL intermediate language to provide formal validation. AIL carries recovered type information and structured control flow, enabling tighter constraints and fewer infeasible paths. This engine plugs directly into MangoDFA, performing targeted symbolic exploration of static candidates to determine path feasibility and generate concrete triggering inputs.
